% Part:sets-functions-relations
% Chapter: size-of-sets
% Section: enumerations-alt

\documentclass[../../../include/open-logic-section]{subfiles}

\begin{document}

\olfileid[fa-IR]{sfr}{siz}{enm-alt}

\olsection{برشماری‌ها و مجموعه‌های \usetoken{S}{enumerable}}

\begin{editorial}
  این بخش برشماری‌ها را، مطابق شیوهٔ رایج در نظریهٔ مجموعه‌ها، به‌صورت
  تناظرهای دوسویی با (قطعه‌های آغازینِ) $\Nat$ تعریف می‌کند. بنابراین
  اندکی با تعریف‌های \olref[enm]{sec} تعارض دارد و همهٔ مثال‌های آنجا
  را تکرار می‌کند. همچنین کمی موجزتر از آن بخش است.
\end{editorial}

می‌توانیم یک مجموعهٔ متناهی را صرفاً با برشمردن
!!{element}sی آن مشخص کنیم. این کار را هنگامی انجام می‌دهیم که مجموعه‌ای را
به‌صورت زیر تعریف می‌کنیم:
\[
  A = \{a_1, a_2, \ldots, a_n\}.
\]
با فرض اینکه !!{element}sی $a_1$، \dots، $a_n$ همگی متمایز باشند،
این تعریف !!a{bijection} میان $A$ و نخستین $n$ عدد طبیعی، یعنی $0$،
\dots، $n-1$، به دست می‌دهد. برعکس، چون هر مجموعهٔ متناهی فقط تعداد
متناهی !!{element} دارد، هر مجموعهٔ متناهی را می‌توان در چنین تناظری
قرار داد. به عبارت دیگر، اگر $A$ متناهی باشد، میان $A$ و
$\{0, \dots, n-1\}$ !!a{bijection} وجود دارد، که در آن $n$ شمار
!!{element}sی~$A$ است.

اگر برخی گونه‌های مجموعهٔ نامتناهی را نیز مجاز بدانیم، آنگاه اجازه
خواهیم داد برخی مجموعه‌های نامتناهی هم برشمرده شوند. می‌توانیم این مطلب
را با گفتن اینکه یک مجموعهٔ نامتناهی به‌وسیلهٔ !!a{bijection} میان آن
و کل~$\Nat$ برشمرده می‌شود، دقیق سازیم.

\begin{defn}[برشماری، مجموعه‌شناختی]
یک \emph{برشماری} از مجموعهٔ $A$، یک !!a{bijection} است که بردش $A$
و دامنه‌اش یا قطعه‌ای آغازین از اعداد طبیعی $\{0,
1, \ldots, n\}$ {یا} کل مجموعهٔ اعداد طبیعی~$\Nat$ است.
\end{defn}

\begin{explain}
این کاربرد واژهٔ \emph{برشماری} مبنایی شهودی دارد. زیرا اینکه بگوییم
مجموعهٔ $A$ را برشمرده‌ایم، یعنی یک !!a{bijection} مانند $f$ وجود دارد
که به ما اجازه می‌دهد عضوهای مجموعهٔ $A$ را یک‌به‌یک بشماریم. عضو
$0$اُم، $f(0)$ است؛ عضو 1اُم، $f(1)$ است؛ \ldots و عضو $n$اُم،
$f(n)$ است\ldots.\footnote{بله، از $0$ می‌شماریم. البته می‌توانستیم
از~$1$ نیز آغاز کنیم. این انتخاب تفاوت مهمی ایجاد نمی‌کرد. فقط لازم
بود~$\Nat$ را با~$\PosInt$ جایگزین کنیم.} مبنای این کاربرد را می‌توان
با افزودن تعریف زیر باز هم روشن‌تر کرد:
\end{explain}

\begin{defn}
  \ollabel{defn:enumerable}
  مجموعهٔ~$A$ !!{enumerable} است اگر و تنها اگر یا $A = \emptyset$
  باشد یا برشماری‌ای از~$A$ وجود داشته باشد. می‌گوییم $A$
  !!{nonenumerable} است اگر و تنها اگر $A$ !!{enumerable} نباشد.
\end{defn}

\begin{explain}
پس یک مجموعه !!{enumerable} است اگر و تنها اگر تهی باشد یا بتوان
به‌کمک یک برشماری !!{element}sی آن را یک‌به‌یک شمرد.
\end{explain}

\begin{ex}
تابعی که اعداد طبیعی را برمی‌شمارد، صرفاً تابع همانی
$\Id{\Nat} \colon \Nat \to \Nat$ است که با $\Id{\Nat}(n) = n$ داده
می‌شود. تابعی که اعداد طبیعیِ \emph{مثبت}، یعنی $\Nat^+ =
\Nat \setminus \{0\}$، را برمی‌شمارد، تابع $g(n) = n + 1$، یعنی تابع
پسین، است.
\end{ex}

\begin{prob}
نشان دهید مجموعهٔ $A$ !!{enumerable} است اگر و تنها اگر یا
$A = \emptyset$ باشد یا !!a{surjection} به‌صورت $f\colon \Nat \to A$
وجود داشته باشد. نشان دهید $A$ !!{enumerable} است اگر و تنها اگر
!!a{injection} به‌صورت $g\colon A \to \Nat$ وجود داشته باشد.
\end{prob}

\begin{ex}
توابع $f\colon \Nat \to \Nat$ و $g \colon \Nat \to \Nat$ که
به‌صورت زیر داده می‌شوند
\begin{align*}
  f(n) & = 2n \text{ و}\\
  g(n) & = 2n+1
\end{align*}
به‌ترتیب اعداد طبیعی زوج و اعداد طبیعی فرد را برمی‌شمارند. اما هیچ‌یک
!!{surjective} نیست؛ پس هیچ‌یک برشماریِ $\Nat$ نیست.
\end{ex}

\begin{prob}
برشماری‌ای از اعداد مربعی $1$، $4$، $9$، $16$، \dots تعریف کنید.
\end{prob}

\begin{ex}
بگذارید $\lceil x \rceil$ تابع \emph{سقف} باشد که $x$ را رو به بالا
تا نزدیک‌ترین عدد صحیح گرد می‌کند. آنگاه تابع $f \colon \Nat \to \Int$
که به‌صورت زیر داده می‌شود:
\[
  f(n) = (-1)^{n} \left\lceil\tfrac{n}{2}\right\rceil
\]
مجموعهٔ اعداد صحیح~$\Int$ را به‌ترتیب زیر برمی‌شمارد:
\[
\begin{array}{c c c c c c c c}
f(0) & f(1) & f(2) & f(3) & f(4) & f(5) & f(6) & \dots \\ \\
\big\lceil \tfrac{0}{2} \big\rceil & -\big\lceil \tfrac{1}{2}\big\rceil &  \big\lceil \tfrac{2}{2} \big\rceil & -\big\lceil \tfrac{3}{2} \big\rceil & \big\lceil \tfrac{4}{2} \big\rceil  & -\big\lceil \tfrac{5}{2}\big\rceil & \big\lceil \tfrac{6}{2} \big\rceil & \dots \\ \\
0 & -1 & 1 & -2 & 2 & -3 & 3& \dots
\end{array}
\]
توجه کنید که $f$ چگونه مقادیر $\Int$ را با «جهیدن» رفت‌وبرگشتی
میان اعداد صحیح مثبت و منفی تولید می‌کند. همچنین می‌توانید $f$ را
به‌صورت موردیِ زیر در نظر بگیرید:
\[
f(n) = \begin{cases}
  \frac{n}{2} & \text{اگر $n$ زوج باشد}\\
  -\frac{n+1}{2} & \text{اگر $n$ فرد باشد}
  \end{cases}
\]
\end{ex}

\begin{prob}
نشان دهید اگر $A$ و $B$ !!{enumerable} باشند، $A \cup B$ نیز چنین است.
\end{prob}

\begin{prob}
با استقرا بر $n$ نشان دهید که اگر $A_1$، $A_2$، \dots، $A_n$ همگی
!!{enumerable} باشند، $A_1 \cup \dots \cup A_n$ نیز چنین است.
\end{prob}

\end{document}
